% Part III: importing tables. Sources: docs/importing.md ("Importing tables: reference", with its
% column reference); docs/databases.md ("Import tables", "Import a known export", "Describe your
% own export in TOML"); docs/administration.md (Databases, Import formats, Imports and sources);
% docs/edf.md (EDF files in imported tables). Checked against medorama/src/tables.rs (columns, synonyms,
% rows, files), imports.rs (batches), linkage.rs (mappings, identity), categories.rs, util.rs
% (times), formats/ (known exports), patient/files.rs (viewer hints, kinds), main.rs (commands),
% web/index.html and web/js/views/sources.js (the page).
\chapter{Importing tables}
\label{chap:importing}

Everything Medorama knows about patients comes from \emph{databases}. In Medorama, a database is
either a table imported into the \emph{catalog} (Medorama's own database of patients,
encounters, entries, users and settings) or a \emph{connected database} that Medorama reads in
place. Most databases are tables exported from an institution's systems: an EHR's patients,
encounters, results and notes, an ECG management system's index, a PACS index, or the bedside
monitors' bed assignments.

This chapter is the reference for importing such tables. It is written for data managers and
for the people who prepare the exports. It states what a table must hold, how its columns are
recognized, how its rows find their patients, how files come in with them, and how a table is
imported, checked and removed. Other ways in are described elsewhere:
\begin{itemize}
\item a database that should stay where it is (an OMOP CDM instance, a research warehouse) can be
  connected instead (\cref{chap:connecting});
\item a CLIF folder is imported table by table (\cref{chap:clif});
\item recordings, ECGs and images come in as files named by such tables (\cref{chap:recordings},
  \cref{chap:ecg-imaging}).
\end{itemize}
The Databases tab, which shows every imported table as it came, is described in
\cref{chap:databases-tab}.

\section{The rule: databases are the only way in}
\label{sec:importing-rule}

A file (a recording, an ECG, a DICOM study, a PDF) enters the catalog only as a row of a table,
and that row names its patient. No file can be added on its own, so no file is ever without its
patient.

Tables can come in \emph{any order}. An item whose patient or encounter no table has named yet
\emph{waits}, and is matched as soon as another table names it. There is one exception: a row
that names only its encounter (no \code{patient\_id}, no \code{mrn}) and holds no item, such as a
row of a table of bed stays by encounter, needs that encounter to be known already, from an
earlier import or an earlier line of the same table (\cref{sec:importing-rows}).

An import can be undone: removing an imported database removes everything it added, and nothing
that another import also added: the links its rows gave recordings go, and a generated patient ID
(\code{GEN-\dots}) that its rows replaced with a patient ID comes back, unless another database
names the patient by that ID. Only administrators import and remove databases, and every import
and every removal is written to the audit log.

An import is all or nothing. Imports run one at a time; one that starts while another runs waits
(its progress says \emph{waiting}). On SQLite, a table is stored a chunk of rows at a time, so
that other changes (a page's edits, the audit log) are written between its chunks; until it ends,
its database is \emph{loading} in the imports list, and what only it adds is not shown on the
pages, in cohorts or in exports. An import that stops part way (an error, a cancel, the server or
the command stopped) is removed as a whole: at once, or when the catalog next starts. A CLIF
folder is one import: a table that fails removes the folder's tables imported before it.

\screenshot{import-flow.pdf}{Importing a database: a table's rows describe patients, encounters and bed stays, carry items, and name files, which become recordings or attachments.}{fig:importing-flow}

\section{Table formats}
\label{sec:importing-formats}

A table is one file with one header row (\cref{tab:importing-formats}). Several files with the
same columns can make one database (\cref{sec:importing-cli}).

\begin{table}[htbp]
\caption{What a table may be.}
\label{tab:importing-formats}
\begin{tabularx}{\linewidth}{@{}P{2.8cm}L@{}}
\toprule
{\tabhead Property} & {\tabhead Rule} \\
\midrule
Kinds & CSV (comma-separated, UTF-8, with or without a byte-order mark), Parquet (in the page
  and by \cmd{medorama import-table}), Excel
  (\file{.xlsx}, \file{.xls}) and OpenDocument (\file{.ods}). Of a spreadsheet, only the first
  sheet is read, and a spreadsheet larger than 100\,MB is refused (it is read whole: save a
  larger table as CSV). The kind is recognized from the file's first bytes, not its name. \\
Encoding and separator & A file whose columns are separated by semicolons or tabs, a UTF-16
  file, and a file that is not UTF-8 (Latin-1, Windows-1252) are refused, with a message that
  says which; they are not guessed at, so that no text is read wrongly. Save the table as CSV in
  UTF-8 (``CSV UTF-8'' in Excel). \\
Header & The first row: one name per column. Names are matched ignoring case, spaces and
  punctuation (\code{Test Name}, \code{test\_name} and \code{TESTNAME} are the same). When two
  columns would be read as the same catalog field (\code{value,value}), the first is; the other
  is kept as imported only, and the report says so. \\
Rows & One row per patient, encounter, bed stay, result, note or file, or several of these at
  once: one row may describe a patient, the patient's encounter and a result. Empty rows are
  skipped. \\
Cells & Text or numbers. Leading and trailing spaces are removed, and an empty cell means ``no
  value''. In a spreadsheet, dates are written as \code{YYYY-MM-DD HH:MM:SS} (or
  \code{YYYY-MM-DD} at midnight) and read in the catalog's time zone, and whole numbers have no
  decimal point. \\
Size & Up to 1\,GiB per upload in the page. The command line has no limit: a CSV or Parquet file
  is read row by row, 5000 rows at a time, so a table of millions of rows needs little memory. A
  spreadsheet is read whole. \\
Kept as imported & Every cell is kept under its column's own name. The Databases tab shows the
  table as it came (\cref{chap:databases-tab}), and removing the database removes everything it
  added. A table of more than 500 columns is kept whole too; criteria on its columns read its
  cells (slower) rather than a table of its numbers. \\
Warnings & A row that cannot be read is named by its line in its file (\code{line 12};
  \code{b.csv, line 12} for a database of several files). \\
\bottomrule
\end{tabularx}
\end{table}

Column names are never changed. A column that the catalog does not recognize is kept and shown
with the table under its own name. Such a column is one of the \emph{database's own columns},
which the import lists and which cohorts can search (\cref{chap:cohorts}). For a table converted
by an import format (\cref{sec:importing-known}), the columns kept are the converted ones.

\section{Columns the catalog recognizes}
\label{sec:importing-columns}

The catalog has a fixed set of columns, the \emph{catalog fields}, in six groups
(\cref{tab:importing-columns}). A table's column is read as one of them when its name is one that
the catalog recognizes, or when it is mapped to it by hand (\cref{sec:importing-mappings}). Each
catalog field takes at most one of the table's columns, and each of the table's columns goes to
at most one catalog field.

Names are compared in lower case, after removing everything but letters and digits.
\Cref{tab:importing-columns} writes the names that way: \code{Patient ID}, \code{patient\_id} and
\code{PATIENTID} are all \code{patientid}. The catalog first looks for each field by its own name
(\code{patient\_id}, \code{bed\_start} and so on), and then by the other names, so that a column
\code{note\_type} meant for \code{name} is not taken by another field first. A column named
\code{unit}, in a table that also has a \code{bed} column and neither a \code{value} nor a
\code{care\_unit} column, is read as the care unit. \Cref{tab:importing-columns} is the complete
list (\code{GROUPS} and \code{SYNONYMS} in \file{medorama/src/tables.rs}).

\begin{xltabular}{\linewidth}{@{}P{3.0cm}P{5.4cm}L@{}}
\caption{The catalog's columns and the names that recognize them (lower case, letters and
digits only).}
\label{tab:importing-columns}\\
\toprule
\tabhead Catalog column & \tabhead Recognized names & \tabhead Meaning \\
\midrule
\endfirsthead
\toprule
\tabhead Catalog column & \tabhead Recognized names & \tabhead Meaning \\
\midrule
\endhead
\bottomrule
\endfoot
\multicolumn{3}{@{}l}{\textbf{Patient}} \\
\code{patient\_id} & \code{patientid}, \code{patient}, \code{subjectid}, \code{subject},
  \code{pid}, \code{personid} & The institution's or the study's patient ID \\
\code{mrn} & \code{mrn}, \code{medicalrecordnumber}, \code{medrecno}, \code{mrno} & The medical
  record number\phimark \\
\code{mrn\_authority} & \code{mrnauthority}, \code{authority}, \code{assigningauthority} & The
  system that issued the MRN, when there are several \\
\code{last\_name} & \code{lastname}, \code{surname}, \code{familyname} & The family
  name\phimark \\
\code{first\_name} & \code{firstname}, \code{givenname} & The given name\phimark \\
\code{birth\_date} & \code{birthdate}, \code{dob}, \code{dateofbirth} & The date of
  birth\phimark \\
\code{sex} & \code{sex}, \code{gender} & \code{F}, \code{M}, or as written \\
\addlinespace
\multicolumn{3}{@{}l}{\textbf{Encounter}} \\
\code{encounter\_id} & \code{encounterid}, \code{encounter}, \code{visitid}, \code{visit},
  \code{admissionid}, \code{hadmid}, \code{stayid} & A visit or a stay \\
\code{csn} & \code{csn}, \code{accountnumber}, \code{account}, \code{fin} & The account or
  financial number\phimark \\
\code{admit\_time} & \code{admittime}, \code{admit}, \code{admissiontime}, \code{admitted},
  \code{admitdate}, \code{admissiondate} & The start of the stay \\
\code{discharge\_time} & \code{dischargetime}, \code{discharge}, \code{dischtime},
  \code{discharged}, \code{dischargedate} & The end of the stay \\
\code{admit\_unit} & \code{admitunit}, \code{admissionunit} & The unit of admission \\
\code{disposition} & \code{disposition}, \code{dischargedisposition} & Where the patient went \\
\addlinespace
\multicolumn{3}{@{}l}{\textbf{Location (a bed stay)}} \\
\code{care\_unit} & \code{careunit}, \code{ward}, \code{location} & The unit. The stays in an
  intensive-care unit make the ICU column \\
\code{bed} & \code{bed}, \code{bedlabel} & The bed \\
\code{bed\_start} & \code{bedstart}, \code{intime}, \code{transferin} & When the patient arrived
  in the bed \\
\code{bed\_end} & \code{bedend}, \code{outtime}, \code{transferout} & When the patient left it \\
\code{recording\_file} & \code{recordingfile}, \code{recording}, \code{recordingname} & A bed
  stay's recording, by its file name or path (ignoring case). It links that recording to the
  encounter directly, instead of by bed and time (\cref{chap:recordings}) \\
\addlinespace
\multicolumn{3}{@{}l}{\textbf{Item}} \\
\code{category} & \code{category}, \code{kind} & The item's column on the patient page
  (\cref{sec:importing-categories}) \\
\code{name} & \code{name}, \code{test}, \code{testname}, \code{item}, \code{parameter},
  \code{label}, \code{drug}, \code{medication}, \code{description}, \code{notetype},
  \code{type}, \code{study} & What was measured, given or written \\
\code{time} & \code{time}, \code{datetime}, \code{charttime}, \code{date}, \code{timestamp},
  \code{resulttime}, \code{collected}, \code{performed}, \code{starttime}, \code{start} & When
  it happened (\cref{sec:importing-times}) \\
\code{end\_time} & \code{endtime}, \code{end}, \code{stop}, \code{stoptime} & When it ended (a
  drug, a support) \\
\code{value} & \code{value}, \code{result}, \code{valuenum}, \code{resultvalue}, \code{dose},
  \code{amount}, \code{measurement} & A number, or a text result \\
\code{unit} & \code{unit}, \code{units}, \code{uom}, \code{valueuom} & The value's unit \\
\code{ref\_low} & \code{reflow}, \code{low}, \code{lower}, \code{normallow}, \code{rangelow} &
  The low end of the normal range \\
\code{ref\_high} & \code{refhigh}, \code{high}, \code{upper}, \code{normalhigh},
  \code{rangehigh} & The high end of the normal range \\
\code{flag} & \code{flag}, \code{abnormal}, \code{abnormalflag} & Abnormal or critical
  (\cref{sec:importing-rows}) \\
\code{code} & \code{code}, \code{icd}, \code{icd10}, \code{cpt}, \code{loinc}, \code{ndc},
  \code{rxnorm} & A code (ICD-10, LOINC and so on) \\
\code{vocabulary} & \code{vocabulary}, \code{codesystem}, \code{system} & The code's
  vocabulary \\
\code{text} & \code{text}, \code{note}, \code{notetext}, \code{report}, \code{reporttext},
  \code{interpretation}, \code{comment}, \code{comments}, \code{body}, \code{impression},
  \code{reaction}, \code{finding}, \code{findings}, \code{conclusion} & A note's or a report's
  text. It is free text, which may name the patient\phimark \\
\code{author} & \code{author}, \code{provider}, \code{clinician}, \code{physician} & Who wrote
  it \\
\code{service} & \code{service}, \code{department}, \code{specialty} & The service \\
\code{domain} & \code{domain} & \code{measurement}, \code{drug}, \code{observation},
  \code{condition}, \code{procedure} or \code{note}; otherwise taken from the category \\
\addlinespace
\multicolumn{3}{@{}l}{\textbf{File}} \\
\code{file} & \code{file}, \code{path}, \code{filepath}, \code{image}, \code{link},
  \code{filename} & A recording or an attachment (\cref{sec:importing-files}) \\
\addlinespace
\multicolumn{3}{@{}l}{\textbf{Other}} \\
\code{source} & \code{source} & The row's source label, which overrides the import's
  \opt{-{}-source} for that row (\cref{sec:importing-batches}) \\
\end{xltabular}

\paragraph{Measurement columns.} A column of the database's own becomes a \emph{measurement
column}, a number that cohorts can search and compare in its unit (\code{QT > 0.4 s}), in two
cases:
\begin{itemize}
\item its name carries a unit in round or square brackets (\code{QT (ms)}, \code{LVEF (\%)},
  \code{Lactate [mmol/L]}); the unit is the text in the last pair of brackets, at most 12
  characters;
\item its name is a known ECG measurement on its own, in its usual unit: \code{Vent. rate},
  \code{Ventricular rate}, \code{Atrial rate}, \code{Heart rate} and \code{HR} are in bpm;
  \code{PR}, \code{QRS}, \code{QT}, \code{QTc}, \code{QTcB}, \code{QTcF} and \code{RR} in ms;
  \code{P axis}, \code{R axis}, \code{QRS axis} and \code{T axis} in degrees.
\end{itemize}
These rules are \code{measure\_column} and the list \code{MEASURES} in
\file{medorama/src/tables.rs}.

\section{How a row is read}
\label{sec:importing-rows}

Each row may carry, in any combination, a patient, an encounter with a bed stay, an item and a
file. Medorama reads a row as follows.
\begin{itemize}
\item \textbf{A patient or an encounter} is read from a row that names a patient
  (\code{patient\_id} or \code{mrn}) and either describes the patient or the encounter (a name,
  a date of birth, the sex, an admission or a discharge, a unit, a disposition, a bed stay, a
  recording file) or holds nothing else (a row of a linkage table). The rules are those of
  \cref{sec:importing-patients}.
\item \textbf{A bed stay by its encounter.} A row that names an encounter but no patient (no
  \code{patient\_id}, no \code{mrn}) and describes the encounter, as in a table of bed stays by
  encounter, belongs to the encounter's patient. When the row holds no item, that encounter must
  already be in the catalog, or on an earlier line of the same table. Otherwise the row is
  refused (\code{line N: no encounter E (import its encounters first)}).
\item \textbf{An item} is read from a row that has a \code{name}, a \code{value}, a \code{text}
  or a \code{file}, and a key: \code{patient\_id}, \code{mrn}, \code{encounter\_id} or
  \code{csn}. A row without a key is refused (\code{no patient\_id, mrn, encounter\_id or csn}).
  The item's name is its \code{name}; otherwise, for a file, the file's name without its
  extension; otherwise the category's label.
\item \textbf{Values.} A \code{value} that is a number (\code{4.1}, \code{-2}, \code{1e3}) is
  stored as a number and can be compared in cohorts; any other value (\code{positive},
  \code{<0.5}) is stored as text. \code{ref\_low} and \code{ref\_high} give the normal range and
  must be numbers. A \code{flag} of \code{HH}, \code{LL}, \code{C}, \code{crit},
  \code{critical}, \code{panic}, \code{AA} or \code{!} marks the value critical (shown in
  orange). A flag of \code{H}, \code{L}, \code{A}, \code{abn}, \code{abnormal}, \code{high},
  \code{low}, \code{pos}, \code{positive}, \code{R}, \code{resistant} or \code{detected} marks
  it abnormal (shown in amber), as does a value outside its range. Flags are compared ignoring
  case.
\item \textbf{Units.} Write the unit as the source system does (\code{mmol/L}, \code{mg/dL},
  \code{bpm}). Cohorts compare values in the unit that a criterion names, converting between
  units where they can.
\item \textbf{Duplicates.} A row identical to an item already stored is not stored again. The
  item is then held by both databases, and stays while either of them does.
\item \textbf{Refused rows.} A row whose time or range cannot be read is refused, with its line
  number and the reason (\code{line 12: not a time: "..."}), and the other rows are imported.
  The report shows at most 50 such warnings.
\end{itemize}

\Cref{tab:importing-smallest} lists the smallest table of each kind. Any other recognized column
may be added.

\begin{table}[htbp]
\caption{The smallest table of each kind.}
\label{tab:importing-smallest}
\begin{tabularx}{\linewidth}{@{}P{3.3cm}P{4.2cm}L@{}}
\toprule
{\tabhead Table} & {\tabhead Required columns} & {\tabhead Usual additions} \\
\midrule
Patients & \code{patient\_id} or \code{mrn} & \code{last\_name}, \code{first\_name},
  \code{birth\_date}, \code{sex}, \code{mrn\_authority} \\
Encounters & \code{encounter\_id}, and \code{patient\_id} or \code{mrn} & \code{admit\_time},
  \code{discharge\_time}, \code{admit\_unit}, \code{disposition}, \code{csn} \\
Bed stays & \code{encounter\_id}, \code{care\_unit}, \code{bed\_start} & \code{bed},
  \code{bed\_end}, \code{recording\_file} \\
Results, vital signs, drugs, diagnoses & \code{patient\_id} or \code{mrn}; \code{name};
  \code{time} & \code{encounter\_id}, \code{value}, \code{unit}, \code{ref\_low},
  \code{ref\_high}, \code{flag}, \code{code}, \code{vocabulary}, \code{category} \\
Notes and reports & \code{patient\_id} or \code{mrn}; \code{time}; \code{text} & \code{name} (the
  note's type), \code{author}, \code{service}, \code{encounter\_id} \\
Medications & \code{patient\_id} or \code{mrn}; \code{name}; \code{time}; the category
  \code{med} & \code{end\_time}, \code{value} and \code{unit} (the dose), \code{code},
  \code{vocabulary} (RxNorm) \\
Conditions, procedures, allergies & \code{patient\_id} or \code{mrn}; \code{name}; the category
  \code{dx}, \code{proc} or \code{allergy} & \code{code}, \code{vocabulary}, \code{time},
  \code{end\_time}; for an allergy, \code{text} (the reaction) \\
Files (recordings, ECGs, imaging, documents) & \code{patient\_id} or \code{mrn}; \code{file} &
  \code{time}, \code{encounter\_id}, \code{category}, measurement columns (\code{QTc (ms)}), the
  bed columns \\
\bottomrule
\end{tabularx}
\end{table}

\section{Categories}
\label{sec:importing-categories}

Each item falls into a \emph{category}, which is its column in the Patient Journey and its group
in the report (\cref{chap:patients}). The categories follow the way a physician reads a chart,
not the systems the data came from: one database may fill several categories, and several
databases one. An item's category is decided in this order:
\begin{enumerate}
\item the category that its row's \code{category} column names;
\item otherwise, the category chosen for the table: \opt{-{}-category} on the command line,
  \ui{Kind of data} or \ui{Items are} in the page, or the category of the table's import format;
\item otherwise, the category that the table's file name or column names suggest: the first word
  of the file name (words of two letters or more), or the first column name, that names a
  category (\file{labs.csv}, \file{ecg\_index.csv}, \file{sleep\_index.csv});
\item otherwise, a category from its domain: a \code{drug} is Medications, a \code{condition}
  Conditions, a \code{procedure} Procedures, a \code{note} Notes and an \code{observation}
  Assessments. A measurement named like a vital sign (\code{HR}, \code{heart rate},
  \code{SpO2}, \code{temp}, \code{SBP}, \code{MAP}, \code{NBP}, \code{blood pressure},
  \code{EtCO2}, \code{CVP}, \code{weight}, \code{height} and so on) is Vitals, and any other is
  Labs.
\end{enumerate}
A category is named by its key, its label or one of its words (\cref{tab:importing-categories}),
in any letter case. The same words find a category in the patient page's search. The list is
\code{CATEGORIES} in \file{medorama/src/categories.rs}.

The \emph{recording categories} are ICU (\code{icu}), Holter (\code{holter}), EEG (\code{eeg})
and Sleep (\code{sleep}). Their files are recordings, indexed in the background and opened in
VitalStream. The files of the other categories are attachments of their item, opened by that
category's viewer (\cref{sec:importing-files}).

\begin{xltabular}{\linewidth}{@{}P{1.4cm}P{2.3cm}L@{}}
\caption{The categories: their keys, labels and the words that name them.}
\label{tab:importing-categories}\\
\toprule
\tabhead Key & \tabhead Label & \tabhead Words \\
\midrule
\endfirsthead
\toprule
\tabhead Key & \tabhead Label & \tabhead Words \\
\midrule
\endhead
\bottomrule
\endfoot
\code{icu} & ICU & icu, intensive care, critical care, monitoring, bedside monitoring, waveforms,
  bedmaster \\
\code{vit} & Vitals & vit, vital, vitals, vital signs, vital sign, flowsheet \\
\code{lab} & Labs & lab, labs, laboratory, chemistry, hematology, coagulation, blood gas, blood
  gases \\
\code{micro} & Microbiology & micro, microbiology, culture, cultures, sensitivity,
  sensitivities \\
\code{path} & Pathology & path, pathology, biopsy, biopsies, cytology, histology \\
\code{ecg} & ECG & ecg, ekg, 12-lead, 12 lead, electrocardiogram, muse \\
\code{holter} & Holter & holter, ambulatory ecg, event monitor, patch monitor \\
\code{eeg} & EEG & eeg, electroencephalography, electroencephalogram, neurophysiology, ieeg,
  seeg \\
\code{sleep} & Sleep & sleep, sleep study, sleep studies, polysomnography, psg \\
\code{echo} & Echo & echo, echocardiography, echocardiogram, tte, tee \\
\code{cath} & Cath lab & cath, cath lab, catheterization, cardiac catheterization, angiography,
  angiogram, coronary angiography, xa, ivus, oct, pci \\
\code{img} & Imaging & imaging, img, radiology, ct, mri, x-ray, xray, ultrasound, dicom, pacs \\
\code{proc} & Procedures & proc, procedure, procedures, surgery, operation, operations, cpt \\
\code{med} & Medications & med, meds, medication, medications, drug, drugs, mar, pharmacy \\
\code{blood} & Transfusions & blood, transfusion, transfusions, blood product, blood products,
  blood bank \\
\code{support} & Life support & support, life support, ventilation, ventilator, respiratory
  support, dialysis, crrt, ecmo \\
\code{dx} & Conditions & dx, condition, conditions, diagnosis, diagnoses, problem, problems,
  problem list, icd \\
\code{obs} & Assessments & obs, observation, observations, assessment, assessments, score,
  scores \\
\code{note} & Notes & note, notes, clinical note, clinical notes, document, documents, report,
  reports \\
\code{allergy} & Allergies & allergy, allergies, intolerance, intolerances (shown with the
  patient, not as a journey column) \\
\end{xltabular}

\section{How rows find their patients}
\label{sec:importing-patients}

The catalog keeps one patient one patient across databases, in the way of a master patient
index: what each database said is kept, the catalog's view is computed from it, and decisions
are kept apart from the data (\cref{chap:administration}). Rows find their patients by the
following rules.
\begin{itemize}
\item \textbf{A patient} is named by \code{patient\_id} (or by an ID that was merged into another
  one), otherwise by an MRN of the same authority (\code{mrn\_authority}). A patient ID is
  compared exactly, as text; MRNs are compared ignoring case and surrounding spaces, and an
  empty authority is the catalog's default one. Give \code{mrn\_authority} whenever MRNs from
  different systems could collide. A row with an MRN but no \code{patient\_id} gets a
  pseudonymous ID from the catalog (\code{GEN-000001} and so on), which stays the same for that
  MRN after any removal and new import, and is never given to another patient. Names, date of
  birth and sex describe a patient; they never identify one on their own.
\item \textbf{Values.} Each field shows the newest batch's value, and removing that batch brings
  back the previous one. A patient may have MRNs of several authorities.
\item \textbf{An encounter} is named by \code{encounter\_id} (otherwise by \code{csn}) and
  belongs to the row's patient. An encounter belongs to one patient: a row that gives it to
  another patient is held back for review.
\item \textbf{An item} (a result, a note, a file) is matched by the first key it has: its
  encounter (\code{encounter\_id}, otherwise \code{csn}), or else its patient (\code{mrn},
  otherwise \code{patient\_id}) and its time, which select the patient's encounter whose stay
  covers that time. An item of a known patient outside every encounter stays with the patient.
  An item that no database names yet waits, and is matched when one does.
\item \textbf{Recordings} also find their encounters by bed and time, by \code{recording\_file},
  or by their BedMaster file names (\code{UNIT\_BED-<time>}; \cref{chap:recordings}).
\item \textbf{Possible duplicates} (two IDs for one person) are listed in \ui{Review}, where an
  administrator with the PHI permission merges them or marks them as different people.
\end{itemize}

\paragraph{Held-back rows.} A row that contradicts the catalog is \emph{held back}: it is kept,
with a warning, for the \ui{Review} tab, and counted as a conflict in the preview
(\cref{tab:importing-held}). The catalog never reassigns a patient on its own on the strength of
an MRN. A merge moves the patient's encounters, MRNs, values, items and sources to the other
patient, and later rows that name the merged ID go to the other patient. A merge cannot be
undone.

\begin{table}[htbp]
\caption{Why a row is held back, and what applying it in Review does.}
\label{tab:importing-held}
\begin{tabularx}{\linewidth}{@{}LL@{}}
\toprule
{\tabhead Problem} & {\tabhead Applying it} \\
\midrule
Its \code{patient\_id} and MRN belong to two different patients & Merges the MRN's patient into
  the ID's patient \\
Its MRN belongs to a patient with another \code{patient\_id} & Records the two IDs as one
  person \\
The patient already has another MRN of the same authority & Adds the MRN \\
Its encounter belongs to another patient & Moves the encounter (with its items) \\
\bottomrule
\end{tabularx}
\end{table}

\section{Times and time zones}
\label{sec:importing-times}

Times are read as \cref{tab:importing-times} describes, in this order. A time without an offset
is local time in the catalog's time zone, the setting \code{timezone} of \file{medorama.toml}
(default \code{America/New\_York}; \cref{chap:settings-reference}). Set it before the first
import, because it decides how every local time is read (\code{parse\_time} in
\file{medorama/src/util.rs}).

\begin{table}[htbp]
\caption{The times a table may hold.}
\label{tab:importing-times}
\begin{tabularx}{\linewidth}{@{}P{6.8cm}L@{}}
\toprule
{\tabhead Written} & {\tabhead Read as} \\
\midrule
\code{1709381400} (only a number greater than $10^8$; a smaller one is not a time, and one of
  $10^{11}$ or more is refused: it counts milliseconds) & Unix seconds \\
\code{2024-03-02T07:10:00-05:00}, \code{2024-03-02T12:10:00Z},
  \code{2024-03-02 12:10:00+00:00} & That moment: the offset or \code{Z} is used, and a space
  may replace the \code{T} \\
\code{2024-03-02 07:10:00}, \code{2024-03-02T07:10:00} (with or without fractions of a
  second), \code{2024-03-02 07:10}, \code{2024-03-02T07:10}, \code{03/02/2024 07:10:00},
  \code{03/02/2024 07:10}, \code{02.03.2024 07:10} & Local time \\
\code{2024-03-02}, \code{03/02/2024}, \code{02.03.2024} & Local midnight \\
Excel and OpenDocument dates & As the spreadsheet stores them, in local time \\
A Parquet timestamp & An instant (\code{isAdjustedToUTC}), or a local time in the catalog's
  time zone when the column says it is not adjusted to UTC \\
\bottomrule
\end{tabularx}
\end{table}

A date with slashes is read in the order the catalog writes its dates (\ui{Admin}
$\rightarrow$ \ui{Dates and times}): day first (\code{03/02/2024} is February 3) when they are
written day/month/year, else month first, the American way (\code{03/02/2024} is March 2).
Choose that setting before importing tables with such dates. A date with dots
(\code{02.03.2024}) is always day first. A year has its four digits: \code{03/02/24} is refused
(\code{write the year with its four digits}), never read as the year 24. A time before 1900 or
after 2200 is refused, and a date of birth before 1880 or in the future. CLIF's times are UTC: a
CLIF table's time written without an offset is read as UTC, not in the catalog's time zone.

A local time that a
daylight-saving change skips (02:30 on the day the clocks go forward) is read with the offset in
force after the change. In New York, 02:30 on March 10, 2024 is read as
\code{2024-03-10T02:30-04:00}, the moment shown as 01:30. Anything else is not a time, and the
row is refused. A date of birth is written \code{YYYY-MM-DD}, \code{MM/DD/YYYY} or
\code{YYYYMMDD}.

\section{Files}
\label{sec:importing-files}

\paragraph{Where.} A \code{file} cell holds a path to a file or a folder. On the command line, a
relative path is read from the table's folder; for a table uploaded in the page, it is read from
the server's working folder, so give absolute paths in tables meant for the page. When the
setting \code{data\_roots} names folders (the \emph{data roots}), every file must lie inside one
of them. An empty \code{data\_roots}, the default, allows any folder; on a server that holds real
data, set it. A file outside the data roots, or one that does not exist, is left out with a
warning (``no such file'', ``is outside the configured data\_roots''), and the rest of the row
is imported. Files are read where they are, never copied, moved or written, and they are shown
only to users who may see the patient.

\paragraph{What they become.} What a file becomes depends on the row's category, or else the
table's (\cref{tab:importing-files}).
\begin{itemize}
\item \textbf{A recording}, when the category is a recording category (ICU, Holter, EEG or
  Sleep) or when neither the row nor the table names one, and the file is one that Medorama can
  read as a recording:
  \begin{itemize}
  \item a WFDB record (its \file{.hea} file; a multi-segment record is taken whole). A WFDB
    header brings in that record only;
  \item a folder. A folder is itself the recording when it is a WFDB folder (one that holds
    WFDB records itself, a \file{.hea} or a \file{RECORDS} file, or an \file{index.csv}, which
    makes it a WFDB collection), read with its subfolders as one session. Any other folder adds
    only the recordings directly inside it (WFDB records, WFDB folders, files that an add-on
    reads; hidden names are skipped), and a subfolder that is not itself a WFDB folder is not
    searched;
  \item an EDF, EDF+, BDF or BDF+ file (\file{.edf} or \file{.bdf}, in any letter case);
  \item a file of a format that an installed add-on reads, such as BedMasterEx \file{.Stp}
    files with stp-adapter (a private add-on, available to institutions under agreement).
  \end{itemize}
  Recordings are added to the catalog, linked to the row's encounter, indexed in the background
  and shown in VitalStream (\cref{chap:recordings}). They take the row's category, and ICU when
  none is given. A recording reaches its patient's page through an encounter: give the row's
  \code{encounter\_id}, or let the recording find its encounter by bed and time. A recording
  that belongs to no encounter is kept in the catalog's list of recordings but is shown on no
  patient's page; a row that names a recording with its patient but no \code{encounter\_id} says
  so in the import report (\code{names no encounter\_id: it is linked only by its bed and time,
  else to no patient}). The encounter a row gives its recording is a link of the import (method
  \code{file}), not a decision by hand: removing the import removes the link, and a decision made
  by hand in \ui{Review} wins over it.
\item \textbf{Nothing}, when such a row names a file that is neither a recording that Medorama
  reads nor a file that a viewer opens (for example a format read by an add-on that is not
  installed). The row is left out with a warning (\code{line N: FILE: not a recording Medorama
  reads. A format read by an add-on needs it installed \dots}), and nothing of it is kept. Once
  the add-on is installed, remove that import and import the table again (\cref{chap:addons}).
\item \textbf{An attachment} of the row's item in every other case, opened by the category's
  viewer. For example, an EDF file whose row has the category \code{ecg} is an attachment, which
  the ECG viewer opens.
\end{itemize}

\begin{table}[htbp]
\caption{What a row's file becomes, by category.}
\label{tab:importing-files}
\begin{tabularx}{\linewidth}{@{}P{2.6cm}LP{2.6cm}@{}}
\toprule
{\tabhead Category} & {\tabhead Files} & {\tabhead Opened in} \\
\midrule
ICU, Holter, EEG, Sleep (or none) & WFDB records, folders of records and collections; EDF, EDF+,
  BDF and BDF+ files; formats read by add-ons, such as BedMasterEx \file{.Stp} with
  stp-adapter & VitalStream, as a recording \\
ECG & A WFDB header with its signal files, GE MUSE XML, a DICOM ECG, an EDF or BDF file, the
  ECG viewer's JSON form, a picture (PNG, JPEG), a PDF & The ECG viewer \\
Echo, Cath lab, Imaging & DICOM files (with or without an extension), a DICOM series' folder,
  videos, pictures (a folder of them, with a frame rate, is one loop), PDF & The PACS viewer \\
Any other & A picture, a PDF, a video, a text file & The entry's detail, on the patient page \\
\bottomrule
\end{tabularx}
\end{table}

\paragraph{In Privacy Mode,} the privacy gate sends only the files it can de-identify or mark:
WFDB headers and their \file{.dat} and \file{.mat} signal files, GE MUSE XML, DICOM, EDF and BDF
files, pictures, videos and PDF. It refuses the rest: JSON ECGs, text files, and signal files
with an extension other than \file{.dat} or \file{.mat} (\code{file\_kind} in
\file{medorama/src/gate.rs}; \cref{sec:protecting-files}, \cref{sec:ecg-imaging-privacy}).

\paragraph{A WFDB record is its header.} A row whose file is a signal file (\file{.dat},
\file{.mat}) is skipped, with a warning that names the header to give in its place
(\cref{sec:wfdb-mat}).

\paragraph{Hints for the viewers.} Columns of the same row tell the viewers what a picture or a
video does not say: its scale (\code{pixel spacing (mm)}, \code{mm per pixel},
\code{pixel size (cm)} and so on) and its frame rate (\code{fps}, \code{frame rate (fps)},
\code{frame rate (Hz)}). The first positive number of each kind counts, and a unit other than
those of \cref{tab:ecg-imaging-hints} is ignored. \Cref{chap:ecg-imaging} describes these hints,
and the databases of ECGs and of DICOM studies that \cmd{muse-py index} and \cmd{dicom-py index}
write from the files themselves.

\section{Import formats}
\label{sec:importing-known}

A known export is recognized by its columns and converted as it is imported. The rows stored are
the converted ones, and the import records its format. On the command line, \opt{-{}-format}
chooses \code{auto} (the default: recognized by its columns), \code{none} (the table as it is),
or a format's key. In the page, \ui{Format} does the same (\ui{Recognize it from its columns} by
default), and a table recognized as a known export is labeled with that format. The built-in
formats are listed in \cref{tab:importing-known}. \cmd{medorama formats} lists them all, with the
format files it could not read, and so does the \ui{Import formats} card of the \ui{Import} tab.

\begin{table}[htbp]
\caption{The built-in import formats.}
\label{tab:importing-known}
\begin{tabularx}{\linewidth}{@{}P{3.4cm}L@{}}
\toprule
{\tabhead Format} & {\tabhead What it is, and what it becomes} \\
\midrule
\code{ge-muse} & GE MUSE ECG metadata, with one row per resting 12-lead ECG and the elements of
  its XML file as columns (\code{PatientDemographics/PatientID},
  \code{RestingECGMeasurements/QTCorrected} and so on). It is recognized by its columns. Each row
  becomes an ECG with its file linked and its patient found by MRN; the measurements become
  columns with their units (\code{QTc (ms)}), and the confirmed and computer statements become
  text. Staff, visit numbers, rooms, beds and the beat list are left out. \\
\code{ge-muse-patient-id} & The same, with the patient found by a patient ID (a research ID, not
  an MRN). It is never recognized on its own: name it with \opt{-{}-format}. \\
\code{clif-\dots} & One format per table of CLIF, the Common Longitudinal ICU data Format, in CSV
  or Parquet (\cref{chap:clif}). \\
\bottomrule
\end{tabularx}
\end{table}

\paragraph{A site's own formats.} A site describes its own exports in TOML files, one per format
(\file{<key>.toml}), in a folder of the setting \code{formats\_dirs} (default
\file{<home>/formats}). No code and no new build are needed. For example:
\begin{lstlisting}[language=TOML]
key = "acme-ecg"                  # --format acme-ecg
label = "ACME ECG export"
category = "ecg"
detect = ["EcgFile", "PatId", "AcqDate"]   # recognized when a table has all these columns
required = ["mrn", "file"]        # rows with these empty are left out (counted)

[[column]]
name = "mrn"                      # a catalog column, or one of the database's own
from = "PatId"
[[column]]
name = "time"
from = "AcqDate"
date = "%m-%d-%Y"                 # read in this form, written as ISO 8601
time_from = "AcqTime"
time = "%H:%M:%S"
[[column]]
name = "QTc (ms)"                 # a measurement: its unit in its name
from = "QTc"
number = true
[[column]]
name = "interpretation"
join = "Statement/{n}/Text"       # numbered columns, in order, joined
\end{lstlisting}
A column that a format does not describe is left out, so a format keeps only what it names. It
should not keep identifying columns other than the patient's (staff, account numbers).
\cmd{medorama formats -{}-check FILE} checks a format file before it is installed. Every key of
such a file is described in \file{medorama/src/formats/declared.rs}, and a worked example is in
the repository's \file{docs/databases.md} (``Describe your own export in TOML''). An export that
a description cannot express needs a converter written in Rust (\cref{sec:readers-add}).

\section{Column mappings}
\label{sec:importing-mappings}

A table whose columns have other names is read through a \emph{mapping}, which pairs each catalog
field with the table's header that holds it.

\paragraph{In the page,} each catalog field has a drop-down list of the table's headers, filled
in from the names that the catalog recognizes. Change any of them. \ui{Save mapping as\dots}
keeps the mapping under a name, and the \ui{saved mapping\dots} list chooses it again for the
next export of the same kind.

\paragraph{On the command line,} \opt{-{}-mapping FILE} gives a mapping in TOML, either under
\code{[columns]} or at the top level of a file that holds nothing else:
\begin{lstlisting}[language=TOML]
# mapping.toml: catalog column = the table's header
[columns]
patient_id = "PAT_ID"
mrn = "Hosp No"
encounter_id = "PAT_ENC_CSN_ID"
name = "Analyte"
value = "Res"
time = "Specimen taken"
\end{lstlisting}
A mapping can also be written in JSON, as one object: \code{\{"mrn": "Hosp No", "name":
"Analyte"\}}. Headers are matched ignoring case, and a key that is not a catalog field is an
error. With a mapping, the other names of \cref{tab:importing-columns} are not looked for. A
header is read as a catalog field only if the mapping names it for that field, or if the header
is that field's own name, ignoring case (\code{mrn}, \code{TIME}, \code{patient\_id}), and the
mapping names no other header for that field. Columns that are not mapped are kept as the
database's own columns, and the import lists them.

\section{Importing in the page}
\label{sec:importing-page}

\screenshot{import.png}{The Import tab's Add a database card, with Import tables chosen.}{fig:importing-page}

Administrators import tables in \ui{Import} $\rightarrow$ \ui{Databases}, in the card \ui{Add a
database}, with the choice \ui{Import tables}:
\begin{enumerate}
\item Choose the \ui{Kind of data} (the category of the tables' items; \ui{Any kind (from the
  files)} by default) and the \ui{Format}.
\item Click \ui{Choose tables on the server\ldots}, open the folder and tick the tables. The
  dialog is the page's own: the server lists its folders, starting at the data roots
  (\code{data\_roots}), and the dialog remembers the last folder. The tables are read where they
  are, row by row, of any size, and a relative path in a \code{file} column is read from the
  table's folder. A table that is only on your own computer is dropped on the drop zone, or
  chosen with \ui{upload from this computer}: up to 1\,GiB per upload, the files of one database
  uploaded together. The files may be CSV, Parquet or spreadsheets. Each file is read at once and shown in a card, with its format if it is a known export, its
  rows, how many of its columns are mapped, \ui{Items are} (its category), the saved mappings,
  its name and its source. Under \ui{Columns}, the card lists the catalog's fields by group, each
  with its drop-down list. Users with the PHI permission also see the file's \ui{First rows}.
  \ui{$\times$} leaves a file out.
\item With several files, choose \ui{Each file a database of its own} or \ui{One database from
  these files (the same columns)}, which needs a name.
\item Click \ui{Preview all} to read every row and count what it would add, keeping nothing
  (\cref{sec:importing-dry-run}).
\item Click \ui{Import all} to import. Files that describe patients and encounters are imported
  first, so that the items of the others find them at once.
\end{enumerate}
Imports run on the server, in the background. You may leave the page: the \ui{Import} tab shows
their progress (rows, share done, time left) wherever you are. \ui{Cancel} stops an import, and
nothing of it is kept, so the catalog stays as it was before. When an import ends, each file's
card reports what was added (new patients, new encounters, bed stays, items, recordings, attached
files), the items already there, the rows held back, the database's own columns, and the
warnings. A table larger than the upload limit is chosen on the server, or imported with the
command line.

\section{Importing on the command line}
\label{sec:importing-cli}

\begin{lstlisting}[style=signature]
medorama import-table FILE... [--format auto|none|KEY] [--mapping FILE] [--category CAT]
                              [--name NAME] [--source LABEL] [--dry-run]
\end{lstlisting}
\begin{params}
\item[FILE...] The table, or several files with the same columns (the same once converted, for a
  known export), which make one database. Each file is read row by row. Several files need
  \opt{-{}-name}.
\item[--format] A known export: \code{auto} (the default), \code{none} or a format's key
  (\cref{sec:importing-known}).
\item[--mapping FILE] A column mapping, in TOML or JSON (\cref{sec:importing-mappings}).
\item[--category CAT] The category of the table's items, when a row names none: a key, a label
  or a word of \cref{tab:importing-categories}.
\item[--name NAME] The batch's name: the database's name in the Databases tab. Default: the
  first file's name.
\item[--source LABEL] The rows' source label. Default: the name. A \code{source} column
  overrides it, row by row.
\item[--dry-run] Check and count only; store nothing.
\end{params}
The following commands show the common uses:
\begin{lstlisting}[style=shell]
medorama import-table encounters.xlsx                    # patients, encounters, bed stays
medorama import-table labs.csv --source labs-2025        # results; the category from the name
medorama import-table ecg_index.csv --category ecg       # ECGs, pictures, PDFs, by their files
medorama import-table icu_index.csv                      # rows linking recordings to encounters
medorama import-table muse/*.csv --name muse_ecg         # several files, one database
medorama import-table export.csv --mapping mapping.toml --source "EHR extract 2024-03"
\end{lstlisting}
Relative file paths are read from the folder of the first file. In a terminal, the import shows
its progress on one line. It then prints one line of counts, such as
(\code{import 12: 865 rows: 0 patient(s) and 0 encounter(s) added, 865 item(s) added (0 already
there, 0 waiting), 0 recording(s) added, 0 file(s) attached, 0 row(s) held back}), and then the
database's own columns and the warnings. A CLIF folder is imported with
\cmd{medorama import-clif} (\cref{chap:clif}). Every command is described in full in
\cref{chap:cli}, and the HTTP requests that the page uses in \cref{chap:http-api}.

\section{Dry runs}
\label{sec:importing-dry-run}

A dry run (\opt{-{}-dry-run}, or the page's \ui{Preview all}) reads and checks every row as the
import does, reports what it would add, match and hold back, and then keeps nothing. It takes
about as long as the import: an ECG list of 500{,}000 rows takes under a minute on SQLite. Run
one before every new kind of table. It shows the columns recognized and those kept as the
database's own, the rows that would be held back, the rows refused, and the items that would
wait.

\section{Examples}
\label{sec:importing-examples}

The following tables use the demo's synthetic patients. A patients table:
\begin{lstlisting}[language=CSV]
MRN,Last Name,First Name,DOB,Sex,Patient ID
SYN0000001,Doe,Jane,1950-03-04,F,P900001
\end{lstlisting}
Encounters, and their bed stays by encounter (imported after the encounters):
\begin{lstlisting}[language=CSV]
Encounter ID,Patient ID,CSN,Admit Time,Discharge Time,Admit Unit
E9100002,P900001,9000000102,2023-12-28 07:40,2024-01-09 13:00,Cardiology
\end{lstlisting}
\begin{lstlisting}[language=CSV]
Encounter ID,Care Unit,Bed,Bed Start,Bed End
E9100002,ICU7,12,2024-01-01 11:00,2024-01-03 11:00
\end{lstlisting}
Laboratory results, matched by MRN and encounter; every column is recognized as it is:
\begin{lstlisting}[language=CSV]
MRN,Encounter ID,Test Name,Result,Units,Ref Low,Ref High,Abnormal Flag,Collected,Code,Code System
SYN0000001,E9100002,Lactate,4.1,mmol/L,0.5,2,H,2024-01-02 06:10,2524-7,LOINC
\end{lstlisting}
An index of ECGs, with their files and measurements:
\begin{lstlisting}[language=CSV]
patient_id,encounter_id,time,QT (ms),QTc (ms),interpretation,file
P900001,E9100002,2024-01-02 08:00,410,452,Sinus tachycardia,ecg/ecg_P900001_20240102.hea
\end{lstlisting}
ICU recordings (their encounters are also found by bed and time, from BedMaster file names such
as \file{ICU7\_12-1704106800.Stp}):
\begin{lstlisting}[language=CSV]
patient_id,encounter_id,file
P900001,E9100002,icu/recording.Stp
\end{lstlisting}
A sleep study: an EDF+ file in the Sleep category, from the demo's \file{sleep\_index.csv} (the
demo's table also names the encounter):
\begin{lstlisting}[language=CSV]
patient_id,encounter_id,time,study,category,file
P900004,E9100010,2024-03-05 22:30,Polysomnography,sleep,../sleep/P900004_psg_20240305.edf
\end{lstlisting}
\file{tools/make-demo-data.py} writes a complete synthetic set of such tables and files, for
trying an import end to end. The installers' \opt{-{}-demo} option loads it into a new catalog as
the cohort DEMO (\cref{chap:quickstart}).

\section{After an import}
\label{sec:importing-after}

\begin{itemize}
\item \textbf{The Databases tab} shows the table as imported, each row with its patient and its
  file (\cref{chap:databases-tab}).
\item \textbf{Review} (administrators) lists the rows held back, the possible duplicate patients
  and the items waiting for their patient or encounter (\cref{chap:administration}).
\item \textbf{Recordings} are indexed by the workers in the background; the Jobs tab shows the
  queue (\cref{chap:recordings}).
\item \textbf{The patients' pages} show each item in its column of the Patient Journey.
\item \textbf{The audit log} holds the import, with its counts.
\end{itemize}

\subsection{Batches and sources}
\label{sec:importing-batches}

Each import is a \emph{batch} with a name and a \emph{source} label. Both default to the table's
file name, and a \code{source} column sets the label row by row. Medorama records what each
batch added. The same recording, patient, encounter or item may come from several batches, and
so may the values: each batch's name, date of birth, admission time and so on are kept apart.
The \ui{Import} tab lists the batches and sources in its \ui{Sources} and \ui{Imports} cards
(\cmd{medorama sources list}, \cmd{medorama imports list}).

\subsection{Removing}
\label{sec:importing-removing}

Removing deletes catalog entries only, and files on disk are never touched
(\cref{tab:importing-removing}). Every removal can be previewed, with \opt{-{}-dry-run} or with
the page's confirmation, which says what would go.

\begin{table}[htbp]
\caption{What each removal takes out.}
\label{tab:importing-removing}
\begin{tabularx}{\linewidth}{@{}P{4.6cm}LP{3.0cm}@{}}
\toprule
{\tabhead Remove} & {\tabhead Takes out} & {\tabhead Leaves} \\
\midrule
A batch (\ui{Imports}, \ui{Remove\dots}; \cmd{imports remove ID}) & What it added, its values,
  its held-back rows & Anything another batch also has; files on disk \\
A source (\ui{Sources}, \ui{Remove\dots}; \cmd{sources remove NAME}) & What every batch added
  under that source & The same \\
The encounters a search finds (\cmd{remove-encounters -{}-search Q}) & Those encounters, with
  their bed stays and items & Recordings (they become unlinked) \\
The recordings a search finds (\cmd{remove -{}-search Q}) & Those recordings (with
  \opt{-{}-purge-results}, their indexes and results too) & Encounters \\
\bottomrule
\end{tabularx}
\end{table}

After a removal, the fields of what stays show the values that the remaining batches gave, links
are computed again (decisions made manually still apply; the links the batch's rows gave its
recordings go), a generated patient ID (\code{GEN-\dots}) that the batch replaced comes back
where no other batch names the patient by the ID it gave, and items whose encounter went away go
back to their patient, or wait. Removing a connected database's batch disconnects it (its
password file in \file{<home>/secrets} goes too). \opt{-{}-purge-results} (on \cmd{imports remove},
\cmd{sources remove} and \cmd{remove}) also deletes the catalog's indexes and results of the
removed recordings, inside \code{results\_dir} and \code{index\_dir} only. The following commands
list the imports, preview a removal, and remove an import and a source:
\begin{lstlisting}[style=shell]
medorama imports list
medorama imports remove 12 --dry-run      # what would go
medorama imports remove 12
medorama sources remove "EHR extract 2024-03"
\end{lstlisting}

\section{Privacy}
\label{sec:importing-privacy}

\begin{phibox}
\begin{itemize}
\item Names, MRNs, account numbers and dates of birth are kept in the catalog's identifier
  tables, and in the imported rows as they came. Users without the PHI permission, and every user
  in Privacy Mode, see placeholders and pseudonyms instead (\cref{chap:protecting}). The first
  rows of a table, in the page, are shown only to users with the PHI permission.
\item Export the identifiers needed for linking and nothing more. A study that needs no names or
  dates of birth should not import them: the catalog links by IDs and MRNs alone.
\item Free text (notes, reports, interpretations) may name patients and staff. By default,
  Privacy Mode shows it with the patient's own names and date of birth replaced, and marked. An
  administrator can withhold it instead (\ui{Admin} $\rightarrow$ \ui{Privacy Mode}: \ui{The
  text of notes and reports}).
\item File paths and names may identify a patient (bedside-monitor systems name files by unit,
  bed and time). They never leave the server in Privacy Mode.
\item Recordings themselves may hold identifiers, such as the patient field of an EDF header.
  The viewers never show them, and in Privacy Mode the gate de-identifies EDF files as they are
  sent (\cref{sec:protecting-edf}).
\item Keep exports for import on an encrypted disk, delete them once imported, and never commit
  them to a repository.
\end{itemize}
\end{phibox}

\section{Checklist}
\label{sec:importing-checklist}

\begin{enumerate}
\item \code{timezone} is set in \file{medorama.toml}, and \code{data\_roots} covers the files'
  folders.
\item The patients table comes first. Any order works, but a first table of patients makes
  matching easy to check. Encounters come before the bed stays that name only their encounter.
\item Each table is a UTF-8 CSV file, a Parquet file or a spreadsheet with one header row. Every
  row names its patient (\code{patient\_id} or \code{mrn}, with \code{mrn\_authority} when MRNs
  come from several systems), and times are in one of the forms of \cref{tab:importing-times}.
\item Categories are given (a \code{category} column, \opt{-{}-category}), or file names and item
  names make them clear. Recordings of EEG and sleep studies have the category \code{eeg} or
  \code{sleep}.
\item Values are numeric where they are numbers, with their units, and codes come with their
  vocabularies.
\item File paths are relative to the table's folder (on the command line) or absolute, inside
  \code{data\_roots}.
\item A dry run shows what would be added, matched and held back.
\item After the import, \ui{Review} shows held-back rows and duplicates, and the Databases tab
  shows the table as imported.
\end{enumerate}
